\documentclass[11pt]{scrartcl}
\usepackage[sexy]{evan} %BLBLBLBLBLBLBLBLBLBLBLB EVANCHEN
\usepackage[T1]{fontenc}
\usepackage[utf8]{inputenc}
\usepackage{graphicx}
\usepackage{amsmath}
\usepackage{tikz}
\usepackage{tikz-cd}
\usepackage{asymptote}
\usepackage{amsthm}
\usepackage{amssymb}
\usepackage{gensymb}
\usepackage{amsfonts}
\usepackage[english]{babel}
\usepackage{quoting}
\quotingsetup{font=big}
\usepackage{booktabs}
\usepackage{caption}
\usepackage{lmodern}
\usepackage{faktor}
\usepackage{bbm} %oer fare gli insiemi indicatori
\usepackage{leftindex} %per gli ideali e gli indici a sinistra
\usepackage{hyperref}
\usepackage{epigraph} 
\usepackage{pgfplots}
\usepackage[export]{adjustbox}

%Pacchetti per i codici
\usepackage{listings}
\usepackage{xcolor}
\usepackage{algpseudocode}

%Analisi e probabilità
%----------------------------------
\newcommand{\dif}{\text{d}}
%%% the pure symbol, see https://tex.stackexchange.com/a/718194/4427
\newcommand{\cind}{\perp\mspace{-9mu}\perp}
%%% define \indep
\NewDocumentCommand{\indep}{e{_}}{%
	\mathrel{%
		\IfNoValueTF{#1}{% no subscript
			\cind
		}{% subscript
			\mathchoice{\mathop{\cind}\limits_{#1}}{\cind_{#1}}{\cind_{#1}}{\cind_{#1}}%
		}%
	}%
}
%-----------------------------------

%Algebra lineare e geometria
%----------------------------------
\newcommand{\tr}[1]{\text{tr}\!\left(#1\right)}
\renewcommand{\rank}[1]{\text{rank}\!\left(#1\right)}
\newcommand{\Immagine}[1]{\text{Im}\left(#1\right)}
\renewcommand{\ker}[1]{\text{Ker}\left(#1\right)}
\renewcommand{\AA}{\mathbb{A}}
\newcommand{\LL}{\mathbb{L}}
\newcommand{\SO}{\text{SO}}
%----------------------------------

%Fisica
%----------------------------------
\newcommand{\UDiMisura}[1]{\,\text{#1}}
\newcommand{\ihat}{\hat{\imath}}
\newcommand{\jhat}{\hat{\jmath}}
\newcommand{\khat}{\hat{k}}
\newcommand{\posit}{\vec{r}}
\newcommand{\veloc}{\dot{\vec{r}}}
\newcommand{\accel}{\ddot{\vec{r}}}
%----------------------------------

%Informatica
%----------------------------------
\definecolor{codegreen}{rgb}{0,0.6,0}
\definecolor{codegray}{rgb}{0.5,0.5,0.5}
\definecolor{codepurple}{rgb}{0.58,0,0.82}
\definecolor{backcolour}{rgb}{0.95,0.95,0.92}

\lstdefinestyle{overleafwiki}{
	backgroundcolor=\color{backcolour},   
	commentstyle=\color{codegreen},
	keywordstyle=\color{magenta},
	numberstyle=\tiny\color{codegray},
	stringstyle=\color{codepurple},
	basicstyle=\ttfamily\footnotesize,
	breakatwhitespace=false,         
	breaklines=true,                 
	captionpos=b,                    
	keepspaces=true,                 
	numbers=left,                    
	numbersep=5pt,                  
	showspaces=false,                
	showstringspaces=false,
	showtabs=false,                  
	tabsize=2
}

\lstset{style=overleafwiki}
\newcommand{\bigo}[1]{\mathcal{O}\!\left(#1\right)}
\newcommand{\bigomega}[1]{\Omega\!\left(#1\right)}
\newcommand{\bigtheta}[1]{\Theta\!\left(#1\right)}
%----------------------------------

%Insiemistica e Algebra
%----------------------------------
\newcommand{\eset}{\varnothing}
\newcommand{\parti}[1]{\mathcal{P}\!\left(#1\right)}
\newcommand{\inj}{\hookrightarrow}
\newcommand{\surj}{\twoheadrightarrow}
\newcommand{\bij}{\hookrightarrow\mathrel{\mspace{-15mu}}\rightarrow}
\newcommand{\ordine}[2]{\text{ord}_{#1}\left(#2\right)}
\newcommand{\normal}{\mathrel{\unlhd}}
\newcommand{\MCD}{\text{MCD}}
\renewcommand{\mcm}{\text{mcm}}
%----------------------------------


\begin{document}
	\title{
		Notes from the Lectures of \\
		Learning from Networks
		}
	\subtitle{Fabio Vandin - UNIPD}
	\date{September 2026 - January 2027}
	
	\maketitle
	
	\tableofcontents
	
	\newpage
	
	\section{Networks analytics}
	
	The reason we are interested in networks is that they are very good at modeling real world problems. The kind of problems we will be concerned about will be posed on graphs the size of the webgraph.
	
	\subsection{Single node level}
	
	The first analytic we can think of is
	\begin{definition}[Node Degree]
		In a graph $(V,E)$ the number of edges incident to a certain node $u$ is called its \textit{degree}, and noted with $\deg u$.
	\end{definition}
	With an adjacency matrix we can compute the degrees for every vertex in $\bigo{n^2}$, but we wish to do better. With and adjancecy list we do a work of
	$$\bigtheta{n+\sum_{v\in V}\deg v}=\bigtheta{n+m}.$$
	Beyond this we have the \textit{node centralities}, a number of scores designed to measure how \textit{central} a node is.
	
	\subsubsection{Closeness}
	
	\begin{definition}[Closeness centrality]
		Let $G=(V,E)$ be a connected undirected graph. We define the \textit{closeness centrality} of a node $v$ as 
		$$c(v):=\frac{n-1}{\sum_{u\in V}d(v,u)}.$$
	\end{definition}
	This is defined to be always $c(v)\in [0,1]$ (except for weigthed graphs, for which we need to renormalize with the sum of all the edges). It's very easy to calculate this analytic for a single node with a Dijkstra or BFS, or a Bellman-Ford for all the graph. The problem is that we actually want to be much faster than quadratic, because the graphs we will be working with are huge. 
	
	\subsubsection{Eppstein-Wang}
	
	We need to approximate the closeness centrality. The main idea is that we want to compute a \textit{campionary average}, that is we need to find the average of distances from all the nodes of the graph to a specific node. So what is natural to do is to compute the distance only for \textit{some} (randomized) subset of the nodes of the graph. Let's be more formal: selecting a random subset $V'$ of vertices from $V$ (not including $v$) we want to say that
	$$c(v)\approx\frac{n-1}{\frac{n}{\abs{V'}}\sum_{v'\in V'}d(v,v')}$$
	considering the simmetric matrix 
	$$D=
	\begin{pmatrix}
		d(1,1) & d(1,2) & \dots & d(1,n) \\
		d(2,1) & d(2,2) & \dots & d(2,n) \\
		\vdots & \vdots & \ddots& \vdots \\
		d(n,1) & d(n,2) & \dots & d(n,n)
	\end{pmatrix}
	$$
	we sample a set $V'$ of vertices of cardinalità $\frac{n}{\varepsilon}$, and we run dijkstra just from the vertices on theese sets. In this way for any node in the graph we have calculated $n/\varepsilon$ distances for it, and we can use them to calculate $c(v)$ as said above. Let's be a bit more formal, we define
	$$\hat{c}(v):=\frac{n-1}{\frac{n}{\abs{V'}}\sum_{v'\in V'}d(v,v')}$$
	the main result we want to show is:
	\begin{proposition}[Central estimate for closeness centrality]
		For a graph $G=(V,E)$, given a random set $V'$ of cardinality $k$ and its associated $\hat{c}$ we want to show that
		$$\mathbb{E}\left[\frac{1}{\hat{c}(v)}\right]=\frac{1}{c(v)}$$
		for any $v$.
	\end{proposition}
	\begin{proof}
		We should first discuss the nature of the random variable $V'$. For $i=1,\dots,k$ we define a succession $X_i$ of iid r.v. (random variables) such that 
		$$X_i:=d(v,v_i)\frac{n}{n-1}$$
		where $v_i$ is the vertex choosen \textbf{uniformly} at random in the $i$-th iteration of the algorithm. Then by definition we have
		$$\frac{1}{\hat{c}(v)}=\frac{1}{k}\sum_{i=1}^{k}X_i$$
		now by linearity of expectation we have
		$$\mathbb{E}\left[\frac{1}{\hat{c}(v)}\right]=\frac{1}{k}\sum_{i=1}^{k}\mathbb{E}\left[X_i\right]$$
		but since the $X_i$ are identically distributed we have $\mathbb{E}[X_i]=\mathbb{E}[X_j]$ and so we just need to show that
		$$\mathbb{E}\left[\frac{1}{\hat{c}(v)}\right]=\frac{1}{k}k\mathbb{E}[X_i]=\mathbb{E}[X_i]$$
		but since the nodes are choosen independently and uniformly with probability $\frac{1}{n}$ we have
		$$\mathbb{E}\left[X_i\right]=\sum_{w\in V}\frac{1}{n}\frac{n}{n-1}d(v,w)=\frac{1}{n-1}\sum_{w\in V}d(v,w)$$
		which is exactly $1/c(v)$, as desired.
	\end{proof}
	But we want more guarantees, so we are going to prove that for $k$ large enough we get close enough approximations of $c(v)$. For this we'll need:
	\begin{proposition}[Hoeffding's Inequality]
		Given $X_i$ a succession of Independent RVs, such that $a_i\le X_i \le b_i$ for every $i$, then by denoting 
		$$\mu := \mathbb{E}\left[\frac{\sum_{i=1}^{k}X_i}{k}\right]$$
		we know that for each $\varepsilon>0$ we have
		$$\mathbb{P}\left[\abs{\frac{\sum_{i=1}^{k}X_i}{k}-\mu}\ge \varepsilon\right]\le 2\exp\left(\frac{-2k^2\varepsilon^2}{\sum_{i=1}^{k}\left(b_i-a_i\right)^2}\right).$$
	\end{proposition}
	\begin{proof}
		This is a consequence of Chernov inequality
	\end{proof}
	The way we are goingo to use this inequality is by introducing:
	\begin{definition}[Diameter]
		For any graph $G=(V,E)$ we call its \textit{diameter} the amount
		$$\Delta(G)=\max_{u,v\in V}d(u,v)$$
		which is a natural generalization for the metric concept of diameter.
	\end{definition}
	\begin{proposition}[Concentration estimate]
		Let $\varepsilon>0$ and $\delta\in (0,1)$ be constants. If $k$ is \textit{big enough}, that is if 
		$$k\ge \frac{1}{2\varepsilon^2}\left(\frac{n}{n-1}\right)^2\log\frac{2n}{\delta}$$
		then we have that
		$$\mathbb{P}\left[\sup_{v\in V}\abs{\frac{1}{\hat{c}(v)}-\frac{1}{c(v)}}\le \varepsilon\Delta(G)\right]\ge 1-\delta$$
		for any $v\in V$.
	\end{proposition}
	\begin{proof}
		We consider again the $X_i$ defined in the central estimate above. We have already proven that
		$$\mu=\frac{1}{c(v)}$$
		so to apply Hoeffding's we just need a bound on the $X_i$s. For this we just consider the diameter, getting
		$$0\le X_i\le \Delta(G)\frac{n}{n-1}.$$
		Now applying Hoeffding it gets quite messy
		\begin{gather*}
			\sum_{i=1}^{k}\left(b_i-a_i\right)\mapsto k\Delta(G)^2\left(\frac{n}{n-1}\right)^2 \\
			\varepsilon \mapsto \varepsilon \Delta(G)
		\end{gather*}
		and so we get
		$$\mathbb{P}\left[\abs{\frac{1}{\hat{c}(v)}-\frac{1}{c(v)}}\ge \varepsilon\Delta(G)\right]\le 2\exp\left(\frac{-2k^2\varepsilon^2\Delta(G)^2}{k\Delta(G)^2\left(\frac{n}{n-1}\right)^2}\right)=2\exp\left(-2k\varepsilon^2\left(\frac{n-1}{n}\right)^2\right)$$
		now we use monotonicity to apply the bound on $k$ and obtain
		$$2\exp\left(-2k\varepsilon^2\left(\frac{n}{n-1}\right)^2\right)\le 2\exp\left(-2\frac{1}{2\varepsilon^2}\left(\frac{n}{n-1}\right)^2\log\frac{2n}{\delta}\varepsilon^2\left(\frac{n}{n-1}\right)^2\right)=\frac{\delta}{n}.$$
		Summaring up we've just proven that
		$$\mathbb{P}\left[\abs{\frac{1}{\hat{c}(v)}-\frac{1}{c(v)}}\ge \varepsilon\Delta(G)\right]\le\frac{\delta}{n}$$
		for any $v\in V$. The event we wish to bound now is
		$$\PP\left[\exists v\in V : \abs{\frac{1}{\hat{c}(v)}-\frac{1}{c(v)}}>\varepsilon\Delta(G)\right]$$
		we can decompose it in 
		$$\left\{\exists v\in V : \abs{\frac{1}{\hat{c}(v)}-\frac{1}{c(v)}}>\varepsilon\Delta(G)\right\}=\bigcup_{v\in V}\left\{\abs{\frac{1}{\hat{c}(v)}-\frac{1}{c(v)}}>\varepsilon\Delta(G)\right\}$$
		and use subadditivity obtaining
		$$\PP\left[\exists v\in V : \abs{\frac{1}{\hat{c}(v)}-\frac{1}{c(v)}}>\varepsilon\Delta(G)\right]\le \sum_{i=1}^{n}\mathbb{P}\left[\abs{\frac{1}{\hat{c}(v)}-\frac{1}{c(v)}}\ge \varepsilon\Delta(G)\right]\le \frac{\delta}{n}n=\delta$$
		taking the complements this gets immediately to the thesis. 
	\end{proof}
	\begin{corollary}
		If $k\in \bigtheta{\frac{\log n}{\varepsilon^2}}$ for a constant $\varepsilon>0$ then we have 
		$$\mathbb{P}\left[\sup_{v\in V}\abs{\frac{1}{\hat{c}(v)}-\frac{1}{c(v)}}\le \varepsilon\Delta(G)\right]\ge 1-\frac{1}{n}$$
		this kind of results are commonly referred to be \textit{true with high probability}.
	\end{corollary}
	As for the complexity, using a dijkstra we get
	$$T(m,n,k)=\bigtheta{k\left(m+n\log n\right)}$$
	and so using $k\in \bigtheta{\frac{\log n}{\varepsilon^2}}$ we get
	$$T(n)=\bigtheta{\frac{m\log n}{\varepsilon^2}}.$$
	
\end{document}
